% hito2_post_ec2_reminder.tex
\documentclass[11pt, a4paper]{article}

\usepackage[T1]{fontenc}
\usepackage[utf8]{inputenc}
\usepackage{tgpagella}
\usepackage{microtype}
\usepackage[spanish, es-noshorthands]{babel}
\usepackage[most]{tcolorbox}
\usepackage[margin=3cm]{geometry}

\definecolor{ucbblue}{RGB}{0, 51, 102}

\begin{document}

\vspace*{2cm}
\begin{center}
    \Large\textbf{\textcolor{ucbblue}{Recordatorio de Avance: Proyecto Hito 2}}\\[0.4cm]
    \large\textit{Siguiente Fase Post-EC2}
\end{center}

\vspace{1cm}
\begin{tcolorbox}[colback=yellow!10, colframe=black, title=\textbf{Transición de Evaluación}, fonttitle=\bfseries]
    El examen EC2 ha evaluado exclusivamente su \textbf{competencia matemática individual}. \\
    
    A partir de este momento, el enfoque de la Unidad II se transfiere totalmente al \textbf{Hito 2}, cuyo propósito es evaluar su capacidad de \textbf{implementación computacional, integración lógica, validación e interpretación}.
\end{tcolorbox}

\vspace{0.5cm}
\textbf{El Hito 2 requiere imperativamente la entrega de:}
\begin{itemize}
    \item Un Jupyter Notebook (\texttt{.ipynb}) que sea completamente reproducible.
    \item Un reporte en formato IEEE en PDF, generado desde compilación LaTeX.
    \item Su asistencia a la defensa técnica grupal con evaluación individual.
\end{itemize}

\vspace{0.5cm}
\begin{tcolorbox}[colback=white, colframe=ucbblue, title=\textbf{Cronograma Provisional}, fonttitle=\bfseries]
    La ventana provisional de desarrollo del proyecto es de aproximadamente \textbf{una semana a partir de la conclusión del EC2}. \\
    
    Las fechas definitivas de entrega del código y defensa técnica serán confirmadas por el docente a la brevedad posible.
\end{tcolorbox}

\vspace{0.5cm}
Se recomienda encarecidamente continuar con la programación del pipeline según las instrucciones establecidas en la \textit{Guía Oficial del Hito 2} y verificar su progreso frente a la Rúbrica Evaluativa previamente entregada.

\end{document}
